\section{Phase~2: Duration-Controlled Speech Synthesis}\label{sec:phase2}
Phase~2 asks whether better phoneme durations give better synthesised speech. We train one acoustic model on the Phase~1 alignments and change only the duration source at test time.

\subsection{System}
\figbox{architecture.png}{Overall system. Phase~1 (forced alignment and duration prediction) supplies the durations that the Phase~2 acoustic model and the pretrained vocoder turn into audio.}{p2arch}
The acoustic model follows FastSpeech~2 \cite{ren2021} without the pitch and energy predictors. Phoneme and stress embeddings with sinusoidal positions pass through four feed-forward Transformer (FFT) blocks (two-head self-attention and a one-dimensional convolutional feed-forward layer). A length regulator repeats each token vector by its duration, and four further FFT blocks produce an 80-bin mel spectrogram. The model has \ptParams{}~M parameters (hidden size 256). A small duration head (two convolutional layers predicting $\log(d{+}1)$) is trained jointly. Training uses teacher forcing: the length regulator receives the forced-alignment durations from Phase~1. Silence is kept as a token, so the durations of an utterance sum exactly to its mel length. The mel recipe (22{,}050~Hz, FFT 1024, hop 256, 80 bins, natural log) is the one the pretrained vocoder expects, and durations are in mel frames (11.6~ms), the same unit as the Phase~1 targets.

\subsection{Training and vocoder}
We use AdamW (peak learning rate $5\!\times\!10^{-4}$, 1{,}000 warm-up steps, inverse-square-root decay), batch size 16 grouped by length, mixed precision and gradient clipping at 1.0. The loss is the masked $L_1$ distance on the normalised mel plus the masked mean squared error of the log-duration. The best checkpoint (step \ptStep{}, validation mel $L_1$ \ptValMel{} in normalised units) was selected on the validation split only. Splits are the frozen Phase~1 splits. Audio is rendered by a pretrained HiFi-GAN \cite{kong2020} (SpeechBrain \cite{ravanelli2021}, trained on LJSpeech), which we did not train. Re-synthesising real clips through our mel recipe and this vocoder gave a mean log-mel $L_1$ of \csMelLone{} against the originals, and the output sounded clean on inspection.

\subsection{Duration-source experiment}
The same trained model synthesised each test utterance with five duration sources: (i)~ground truth (upper bound), (ii)~the Phase~1 BiLSTM, (iii)~the phone-mean baseline, (iv)~the global mean, and (v)~the acoustic model's own duration head. Phase~1 predicts speech phones only, so silence tokens keep their ground-truth length in every condition and only speech-phone durations differ. Because different durations give different lengths, we compare the synthesised mel $\hat M$ with the real mel $M$ after dynamic time warping \cite{sakoe1978}:
\begin{equation}
D(\hat M, M)=\frac{1}{|\pi^\star|\,F}\sum_{(i,j)\in\pi^\star}\lVert \hat m_i-m_j\rVert_1,\quad F=80,
\end{equation}
where $\pi^\star$ is the minimum-cost warping path. We also report the absolute error of total length. Intervals are 95\% bootstrap CIs over \ptEvalN{} test utterances, and conditions are compared by paired bootstrap differences on the same utterances.

\subsection{Results}
\ptFindings{}

\safeinput{tables/phase2_objective.tex}
\safeinput{tables/phase2_paired.tex}
\figbox{p2_objective.png}{Duration-source ablation: DTW mel distance (left) and absolute total-length error (right), with 95\% bootstrap intervals.}{p2obj}
\figbox{p2_curves.png}{Training and validation curves of the acoustic model.}{p2curves}

\subsection{Subjective evaluation}
\listeningParagraph{}
\IfFileExists{tables/phase2_mos.tex}{\input{tables/phase2_mos.tex}}{}

\subsection{Limitations of Phase~2}
\begin{itemize}
\item The DTW mel distance is a proxy for quality: it removes timing differences, so it measures spectral similarity given the durations, not naturalness.
\item The acoustic model is small and has no pitch or energy predictors; its own mel error dominates the remaining distance, which limits how much better durations can help.
\item The vocoder is pretrained, the data are one speaker of read English, and the experiment uses \ptEvalN{} test utterances.
\item The text demonstration uses the model's own duration head; driving it with the Phase~1 BiLSTM on free text requires exposing the Phase~1 feature builder as a function (future work).
\end{itemize}
